% ICNAAO 2026 abstract template
% Replace the sample title, author information, abstract, keywords, and references.
% Move the "Presenting author" label to the author who will give the talk.
% Compile twice with: pdflatex abstract-template.tex
\documentclass[11pt,a4paper]{article}

\usepackage[T1]{fontenc}
\usepackage[utf8]{inputenc}
\usepackage{amsmath,amssymb,amsfonts,amsthm}
\usepackage[margin=1in]{geometry}
\usepackage[hidelinks]{hyperref}

\pagestyle{empty}

\title{Title of Your Contribution}
% Define each distinct affiliation once with \thanks{...}.
% Reuse its mark with \footnotemark[1] (or [2], etc., in order of \thanks).
% For an author at a different institution, replace \footnotemark[1]
% with \thanks{Other Department, Other University, City, Country.}.
% Keep author-specific details outside the shared affiliation footnote.
\author{%
  First Author\thanks{Department of Mathematics, First University, City, Country.}\\
  \small\href{mailto:first@example.com}{\texttt{first@example.com}}
  \and Second Author\footnotemark[1]\\
  \small Presenting author\\
  \small\href{mailto:second@example.com}{\texttt{second@example.com}}
  \and Third Author\thanks{Department of Mathematics, Second University, City, Country.}\\
  \small\href{mailto:third@example.com}{\texttt{third@example.com}}
}
\date{}

\begin{document}
\maketitle
\thispagestyle{empty}

\begin{abstract}
Replace this text with a concise description of your research problem,
methods, main results, and significance. 
For example, consider the optimization problem
\[
    \min_{x \in C} f(x),
\]
where $C \subseteq \mathbb{R}^{n}$ is a nonempty closed convex set and
$f\colon \mathbb{R}^{n} \to \mathbb{R}$ is continuously differentiable.
Describe your assumptions, the proposed method, and the principal
conclusions here. If needed, cite relevant literature using
\cite{sample-article,sample-book}.
\end{abstract}

\noindent\textbf{Keywords:} First keyword; second keyword; third keyword.

% Optional: uncomment and supply the relevant classification codes.
% \medskip
% \noindent\textbf{Mathematics Subject Classification:} Primary XXXXX; Secondary XXXXX.

% Optional references: replace these illustrative placeholders with actual
% sources, or remove this section and the sample citation if not needed.
% Use \cite{key} in the abstract to refer to a matching \bibitem{key} below.
\begin{thebibliography}{9}
\bibitem{sample-article}
A.~Author and B.~Author,
Title of the article,
\emph{Journal Name} \textbf{Volume} (Year), pages.

\bibitem{sample-book}
C.~Author,
\emph{Title of the Book},
Publisher, Year.
\end{thebibliography}

\end{document}
